1.
\bf{Introduction générale}
\bf{Notion d'état d'un système physique :}
-Cas classique et cas quantique
-Évolution d'un système
\bf{Notion d'"espace des états" :}
Principe de combinaison
Structure de l'espace des états : espace vectoriel
Corps des scalaires : nombres complexes
\bf{Notation des vecteurs d'états : "kets"}
\bf{ATTENTION :} mise en garde essentielle à propos des combinaisons linéaires d'états

2.
Évolution d'un système physique
Opérateur d'évolution
Comment distinguer les états physiques ?
- États à valeur bien définie d'une grandeu donnée
- États sans valeur bien définie
- Grandeurs physiques et combinaisons linéaires
Problématique générale de la mesure :
-interaction système/instrument
interprétation de Copenhague, interprétation d'Everett

3.
Problématique générale de la mesure (fin)
- Indéterminisme, aléatoire, probabilités
- Règle de Born
- Impossibilité de prendre connaissance d'un état initialement inconnu
Structure hermitienne de l'espace des états
- produit scalaire sur un espace vectoriel réel
Attention : la dernière demi-heure est sans son, en raison d'une défaillance de l'alimentation de système audio... Toute mes excuses !

4.
\bf{}
\bf{}
\bf{}

5.
Adjoint d'un opérateur
Opérateurs auto-adjoints
Construction de l'espace des états d'un système physique :
- base orthonormée de vecteurs à valeur bien définie d'une grandeur physique
-"coordonnées" d'un vecteur d'état
- relation de fermeture, projecteur orthogonal
- mesure d'une grandeur physique

6.
=== Fin du chapitre 1
Espace des états d'un système quantique
Grandeurs physiques complémentaires
Notion d'observable physique
"Compatibilité" entre observables

=== Début du chapitre 2
Systèmes quantiques élémentaires
Notion de qubit

7.
Systèmes quantiques élémentaires : qubit
- représentation d'un qubit sur la sphère de Bloch
- représentation d'un ket en vecteur colonne
- représentation matricielle d'un opérateur
- observable sur un qubit (opérateur auto-adjoint)
- matrices de Pauli, valeurs et vecteurs propres
Expérience de Stern et Gerlach (1922)
- Rappel : moment magnétique
- Dispositif de Stern et Gerlach, analyse classique

8.
Expérience de Stern et Gerlach (1922)
-Description du dispositif et analyse classique : mesure d'une composant particulière du moment magnétique
-Résultats de l'expérience
-Résultat de mesure effectuées en cascade
Description dans le cadre du formalisme quantique
- Espace des états et base quelconque de l'espace
- Correspondance entre états : combinaison linéaires

9.
Opérateur vectoriel "moment magnétique"
Lien avec les matrices de Pauli
Composante du moment magnétique selon un axe quelconque : vecteurs propres, valeurs propres
Orientation dans l'espace géométrique 3D :
- lien avec la "sphère de Bloch"
- valeurs moyennes de scomposantes
Notion de spin : moment angulaire intrinsèque
Spin 1/2

10.
\bf{Spin :}
- opérateur vectorie (moment cinétique intrinsèque)
- relation de comutation
- nombre quantique de spin
- espace vectoriel associé et valeurs propres
\bf{Combinaison de deux systèmes quantiques :}
- Espace des états
- Produit tensoriel
- Exemple de deux qubits
- États séparables et états intriqués

11.
\bf{Combinaison de deux systèmes quantique}
- Espace des états : produit tensoriel
- Action linéaire dans l'espace produit tensoriel : produit tensoriel de deux opérateurs linéaires
- "Promotion" d'un opérateur agissant sur un seul sous-système
- Exemple de deux qubits
- Mesure individuelles sur un qubit intriqué

12.
\bf{Combinaison de deux systèmes quantiques}
- Situation de "paradoxe" E.P.R. (Einstein, Podolsky, Rosen)
- Discussion de l'"influence instantanée à distance" (ce qui n'a pas le moindre sens en Physique !)
- Problème central dans l'interprétation de Copenhague
- Description de la situation dans l'interprétation d'Everett
\bf{CHAPITRE 3 : Évolution dans le temps des systèmes}
- Opérateur évolution : linéarité, invarience par translation
- Unitarité (persistance de la distinction entre états)

13
\bf{CHAPITRE 3 : Évolution dans le temps des systèmes}
- Opérateur évolution : unitarité
- Évolution infinitésimale
- Opérateur Hamiltonien
- Équation de Schrödinger
- États stationnaires
- Conservation de l'énergie

14.
\bf{Chap. 3 : Évolution dans le temps des systèmes}
- Variation d'une grandeur physique quelconque (rôle du commutateur avec le hamiltonien)
- Exemple d'un spin 1/2 dans un champ magnétique
\bf{Chap. 4 : Espace des états de position}
- Introduction

15
\bf{Chap. 4 : Espace des états de position}
- États de position : existe-t-il des états correspondant à une position bien définie ?
- Peut-on donner un sens à une combinaison linéaire" continue ?
- Notion de "bra généralisé" : des formes linéaires continue aux distributions
- "Ket généralisé"
- Cas des bras et des kets généralisés "impulsion"

16.
\bf{Espace des états de position}
Bilan de l'étude :
- espace des fonctions de carré sommable : L2(R)
- bras et kets généralisés
- bases généralisées : \{<p|\} et \{<x|\}
- <x| : évaluation de la fonction en x
- <p| : évaluation de la transformée de Fourier en p
- relation de fermeture
Représentations "position" et "impulsion"
Fonction d'onde et signification physique

17.
\bf{Chap. 4 : Espace des états de position}
\bf{III - Oopérateurs dans l'espace des états :}
Opérateur position (à 1D et à 3D) :
- action sur les vecteurs d'états, en représentation "position" et en représentation "impulsion"
Opérateurs "translation" : unitarité, composition
Générateur des translations :
- opérateur auto-adjoint, P : opérateur "impulsion"
- action en représentations "position" et "impulsion"

18.
\bf{Chap. 4 : Espace des états de position}
\bf{Résumé :} opérateurs position et impulsion, action de ces opérateurs en représentation position et en représentation impulsion
Relation de commutation (r et p)
Écart type des mesures d'une grandeur physique quelconque
Relation d'indétermination (Heisenberg)
Observables compatibles / incompatibles
Cas particulier : position vs. impulsion
Opérateurs de rotation

19.
\bf{Chap. 4 : Espace des états de position (fin)}
Opérateurs générateurs des rotations
Moment cinétique
Relation de commutation fondamentales
\bf{Chap. 5 : Equation de Schrödinger et "quantification" des systèmes classiques}
"Particule" libre dans l'espace
Hamiltonien libre
États stationnaires (représentations "r" et "p")

20.
\bf{Chap. 5 : Equation de Schrödinger et "quantification" des systèmes classiques}
Cas d'un "particule" libre :
- États stationnaires en représentations "r" et "p"
- Équation de Schrödinger "indépendante du temps"
- Évolution de la moyenne de la position
- Évolution de la moyenne de l'impulsion
Cas d'une "particule" soumise à un potentiel :
- Hamiltonien et équation de Schrödinger
- Évolution moyenne de la position et de l'impulsion
- "Force" apparente

\bf{Chap. 5 : "Quantification" des systèmes classiques}
Conservation locale de la probabilité :
- interprétation probabiliste de la fonction d'onde : densité de probabilité de "présence"
- équation de continuité
- vecteur "densité de courant de probabilité"
Exemple d'un puits de potentiel à une dimension :
- solutions stationnaires

22.
\bf{Chap. 5 : "Quantification" des systèmes classiques}
Puits de potentiel fini et puits de pottentiel infini :
- solutions stationnaires de l'équation de Schrödinger
- bilan des paramètres libres et des contraintes
- énergies possibles discrètes : "quantification"
- continuum d'énergies possibles
Marche de potentiel
- "transmission" et "réflexion" en amplitude
- "effet tunnel"
L'oscillateur harmonique quantique (introduction)

23.
\bf{Chap. 5 : "Quantification" des systèmes classiques}
L'oscillateur harmonique quantique
- hamiltonien
- recherche des états stationnaires
- opérateurs d'échelle : a et a+
- opérateur "nombre de quanta" : N
- relations de commutation
- valeurs propres de N et H
- vecteurs propres associés
- base de Fock et polynômes de Hermite
